\BourbakiExercisesSection

\begin{enumerate}
\def\labelenumi{\arabic{enumi}.}
\item
  设 N 为 Z-模 Z/4Z，M 为 N 的子模 2Z/4Z，且 \(j:M\to N\) 为典范单射。证明 \(\mathsf{T}(j):\mathsf{T}(\mathsf{M})\to\mathsf{T}(\mathsf{N})\) 不是单射的。
\item
  设 I 和 J 是两个有限集合， \(I_{0} = I \cup \{\alpha\}\) , \(J_{0} = J \cup \{\beta\}\) ，其中 \(\alpha\) 与 \(\beta\) 分别是不属于 I 与 J 的元素，并设 \(I_{0}\) 与 \(J_{0}\) 不相交。设 E 是交换域 K 上有限秩的向量空间，z 是 \(\mathbf{T}_{J}^{\mathrm{I}}(\mathbf{E})\) 的一个元素。对所有 \(i \in I\) ，设 \(W_{i}'\) 是 \(E^{*}\) 的由元素 \(x^{*} \in E^{*}\) （满足 \(c_{\beta}^{i}(z \otimes x^{*}) = 0\) ）组成的子空间（\(c_{\beta}^{i}\) 是指标 \(i \in I\) 与指标 \(\beta \in J_{0}\) 在 \(\mathbf{T}_{J_{0}}^{\mathrm{I}}(\mathbf{E})\) 中的缩并），并设 \(V_{i}\) 是 E 中与 \(W_{i}'\) 正交的子空间。对所有 \(j \in J\) ，设 \(W_{j}\) 是 E 的由元素 \(x \in E\) （满足 \(c_{j}^{\alpha}(z \otimes x) = 0\) ）组成的子空间（\(c_{j}^{\alpha}\) 是指标 \(\alpha \in I_{0}\) 与指标 \(j \in J\) 在 \(T_{J}^{I_{0}}(E)\) 中的缩并），并设 \(V_{j}'\) 是 \(E^{*}\) 中与 \(W_{j}\) 正交的子空间。证明 z 属于张量积 \(\left(\bigotimes_{i \in I} V_{i}\right) \otimes \left(\bigotimes_{j \in J} V_{j}'\right)\) ，它典范地等同于 \(T_{J}^{I}(E)\) 的一个子空间。此外，若 \((\mathbf{U}_{i})_{i \in I}\) 是 E 的一族子空间， \((\mathbf{U}_{j}')_{j \in J}\) 是 \(E^{*}\) 的一族子空间，使得 z 属于张量积 \(\left(\bigotimes_{i \in I} U_{i}\right) \otimes \left(\bigotimes_{j \in J} U_{j}'\right)\) ，它典范地等同于 \(T_{J}^{I}(E)\) 的一个子空间，则 \(V_{i} \subset U_{i}\) 且 \(V_{j}' \subset U_{j}'\) 对所有 \(i \in I\) 和 \(j \in J\) 成立（参见 II，§7，第 8 号，命题 17）。
\item
  设 M 是有限生成的投射 A-模。对 M 的每个自同态 u，用 \(\tilde{u}\) 表示与之对应的张量 \(\theta_{\mathbf{M}}^{-1}(u) \in \mathbf{M}^{*} \otimes_{\mathbf{A}} \mathbf{M}\) 。
\end{enumerate}

\begin{enumerate}
\def\labelenumi{(\alph{enumi})}
\item
  对每个元素 \(x \in \mathbf{M}\) ，证明 \(u(x) = c_2^1(x \otimes \tilde{u})\) ，其中 \(x \otimes \tilde{u}\) 属于 \(\mathbf{M} \otimes \mathbf{M}^* \otimes \mathbf{M} = \mathsf{T}_{\langle 2\rangle}^{(1,3)}(\mathbf{M})\) 。
\item
  设 \(u, v\) 是 \(M\) 的两个自同态且 \(w = u \circ v\) ；证明 \(\tilde{w} = c_3^2(\tilde{v} \otimes \tilde{u})\) ，其中 \(\tilde{v} \otimes \tilde{u}\) 属于 \(T_{(1,3)}^{(2,4)}(M)\) 。
\item
  对每个自同构 \(s\) （作用于 \(\mathbf{M}\)）与每个张量 \(z \in \mathbf{T}_{\mathbf{J}}^{\mathbf{I}}(\mathbf{M})\) ，张量 \(s, z \in \mathbf{T}_{\mathbf{J}}^{\mathbf{I}}(\mathbf{M})\) 由如下条件定义：若 \(z = \bigotimes_{k \in \mathbf{I} \cup \mathbf{J}} z_k\) ，其中 \(z_k \in \mathbf{M}\) （若 \(k \in \mathbf{I}\)）而 \(z_k \in \mathbf{M}^*\) （若 \(k \in \mathbf{J}\)），则 \(s, z = \bigotimes_{k \in \mathbf{I} \cup \mathbf{J}} z_k'\) ，其中 \(z_k' = s(z_k)\) （若 \(k \in \mathbf{I}\)），\(z_{k}^{\prime}=^{t}s^{-1}(z_{k})\) （若 \(k\in J\)）。证明群 \(\mathbf{GL}(\mathbf{M})\) 以法则 \((s,z)\mapsto s.z\) 作用于 A-模 \(T_{J}^{I}(M)\) 上，诸映射 \(z\mapsto s.z\) 是该 A-模的自同构。此外，对每对有序指标 \(i\in I\) , \(j\in J\) ,
\end{enumerate}

\[
c _ {j} ^ {i} (s. z) = s. c _ {j} ^ {i} (z).
\]

若 M 是投射的且有限生成的，则在习题 2 的记号下证明 \(s. \tilde{u} = (s \circ u \circ s^{-1})^{\sim}\) 。

\begin{enumerate}
\def\labelenumi{\arabic{enumi}.}
\setcounter{enumi}{3}
\tightlist
\item
  设 M 是有限生成的投射 A-模；对每个双线性映射 \(u\) （从 \(\mathbf{M} \times \mathbf{M}\) 到 \(\mathbf{M}\) 中），用 \(\tilde{u}\) 表示张量 \(\theta_{\mathbf{M}}^{-1}(u)\) （属于 \(\mathsf{T}_{(1,2)}^{(3)}(\mathbf{M})\)）。证明：\(\mathbf{M}\) 上由 \(u\) 定义的代数结构是结合的，当且仅当张量 \(c_4^3 (\tilde{u} \otimes \tilde{u})\) 与 \(c_2^6 (\tilde{u} \otimes \tilde{u})\) 在从 \(\mathsf{T}_{(1,2,3)}^{(4)}(\mathbf{M})\) 到 \(\mathsf{T}_{(1,3,4)}^{(2)}(\mathbf{M})\) 上的典范结合同构下互相对应。
\end{enumerate}

¶ 5. 设 E 是交换域 K 上的向量空间，并设 T(E) 是 E 的张量代数。

\begin{enumerate}
\def\labelenumi{(\alph{enumi})}
\item
  证明 \(\mathsf{T}(\mathbf{E})\) 不含零因子，且 \(\mathsf{T}(\mathbf{E})\) 中仅有的可逆元素是标量 \(\neq 0\) 。若 \(\dim (\mathbf{E}) > 1\) ，则 \(\mathsf{T}(\mathbf{E})\) 是非交换 K-代数。
\item
  设 u、v 是 T(E) 的两个元素。证明：若存在 T(E) 的两个元素 a、b 使得 \(ua = vb \neq 0\) ，则 u、v 中之一是另一个的右倍元。（先考虑 u 与 v 都是齐次的情形；把它化归为 E 是有限维的情形，并写出 \(u = e_{1}u_{1} + \cdots + e_{n}u_{n}\) , \(v = e_{1}v_{1} + \cdots + e_{n}v_{n}\) ，其中 \((e_{k})\) 是 E 的一个基。再对 ua 的齐次分量的最小次数作归纳得出结论。）由此推出：若 \(\dim(\mathrm{E}) > 1\) ，环 T(E) 没有左分式域（I，§ 9，习题 15）。
\item
  证明：两个非零元素 \(u, v\) （属于 \(\mathsf{T}(\mathbf{E})\)）可交换，当且仅当存在向量 \(x \neq 0\) （属于 \(\mathbf{E}\)），使得 \(u\) 与 \(v\) 属于 \(\mathsf{T}(\mathbf{E})\) 的由 \(x\) 生成的子代数（利用 (b)）。由此推出：若 \(\dim(\mathbf{E}) > 1\) ，则 \(\mathsf{T}(\mathbf{E})\) 的中心等于 \(\mathbf{K}\) 。
\end{enumerate}

¶ 6. 设 K 是交换域，E、F 是两个 K-代数，各带单位元，且该单位元与 K 的单位元 l 等同。考虑 K-代数 T(E ⊕ F)，其中 E ⊕ F 被看作 K 上的向量空间；在典范单射下，E（相应地 F）等同于 T(E ⊕ F) 的一个向量子空间。设 \(\mathfrak{J}\) 是 T(E ⊕ F) 中由元素 \(x_1 \otimes x_2 - (x_1 x_2)\) 与 \(y_1 \otimes y_2 - (y_1 y_2)\) 生成的双边理想，其中 \(x_i \in E\) , \(y_i \in F\) , \(i = 1, 2\) ；设 E * F 表示商代数 T(E ⊕ F)/\(\mathfrak{J}\)；它称为 K-代数 E 与 F 的自由积；设 \(\phi\) 是 T(E ⊕ F) 到 E * F 上的典范同态，\(\alpha\) 与 \(\beta\) 是 \(\phi\) 在 E 与 F 上的限制；\(\alpha\) 与 \(\beta\) 是 K-代数同态。

\begin{enumerate}
\def\labelenumi{(\alph{enumi})}
\item
  证明：对每一对 K-同态 \(u: \mathbf{E} \to \mathbf{G}\) , \(v: \mathbf{F} \to \mathbf{G}\) （取值在 K-代数 \(\mathbf{G}\) 中），存在唯一的 K-同态 \(w: \mathbf{E} * \mathbf{F} \to \mathbf{G}\) ，使得 \(u = w \circ \alpha\) 且 \(v = w \circ \beta\) （这就证明了该名称的合理性）。
\item
  设 \(R_{n}\) 是 \(\mathsf{T}(\mathbf{E}\oplus\mathbf{F})\) 中诸 \(\mathsf{T}^{k}(\mathbf{E}\oplus\mathbf{F})\) 对 \(0\leqslant k\leqslant n\) 之和，并记 \(J_{n}=J\cap R_{n}\) 。证明：若考虑由 l 与族 \((e_{\lambda})_{\lambda\in\mathbb{L}}\) 组成的 E 的基，以及由 l 与族 \((f_{\mu})_{\mu\in\mathbb{M}}\) 组成的 F 的基，则 \(J_{n}+R_{n-1}\) 在 \(R_{n}\) 中的一个补子空间可如下得到：在 \(R_{n}\) 中考虑以形如 \(z_{1}\otimes z_{2}\otimes\cdots\otimes z_{n}\) 的元素为基的子空间，其中或者 \(z_{2j+1}\) 等于某个 \(e_{\lambda}\) （\(2j+1\leqslant n\)）且 \(z_{2j}\) 等于某个 \(f_{\mu}\) （\(2j\leqslant n\)），或者 \(z_{2j+1}\) 等于某个 \(f_{\mu}\) （\(2j+1\leqslant n\)）且 \(z_{2j}\) 等于某个 \(e_{\lambda}\) （\(2j\leqslant n\)）（对 n 作归纳）。
\item
  设 \(\mathbf{P}_n = \phi (\mathbf{R}_n)\subset \mathbf{E}*\mathbf{F}\) ；于是 \(\mathbf{P}_n\subset \mathbf{P}_{n + 1},\mathbf{P}_0 = \mathbf{K},\mathbf{P}_h\mathbf{P}_k\subset \mathbf{P}_{h + k}\) ，且 \(\mathbf{E}*\mathbf{F}\) 是诸 \(\mathbf{P}_n\) 之并。证明存在典范的向量空间同构，对 \(n\geqslant 1\)
\end{enumerate}

\[
\mathrm{P} _ {n} / \mathrm{P} _ {n - 1} \rightarrow \left(\mathrm{G} _ {1} \otimes \mathrm{G} _ {2} \otimes \dots \otimes \mathrm{G} _ {n}\right) \oplus \left(\mathrm{G} _ {2} \otimes \mathrm{G} _ {3} \otimes \dots \otimes \mathrm{G} _ {n + 1}\right)
\]

其中记 \(G_{k} = E/K\) （k 为奇数），\(G_{k} = F/K\) （k 为偶数）。特别地，同态 \(\alpha\) 与 \(\beta\) 是单射，这使我们可以把 E 与 F 等同于 E * F 的子 K-代数。用 \(P_{n}^{\prime}\) （相应地 \(P_{n}^{\prime\prime}\) ）表示 \(G_{1} \otimes G_{2} \otimes \cdots \otimes G_{n}\) （相应地 \(G_{2} \otimes G_{3} \otimes \cdots \otimes G_{n+1}\) ）在 P 中的逆像；元素 \(z \in E * F\) 的高度（记作 \(h(z)\) ）是使 \(z \in P_{n}\) 成立的最小整数；若 \(h(z) = n\) 且 \(z \in P_{n}^{\prime}\) （相应地 \(z \in P_{n}^{\prime\prime}\) ），则称 z 为纯奇的（相应地，纯偶的）。于是 \(P_{n} = P_{n}^{\prime} + P_{n}^{\prime\prime}\) 且 \(P_{n}^{\prime} \cap P_{n}^{\prime\prime} = P_{n-1}\) 对 \(n \geqslant 1\) 成立。

\begin{enumerate}
\def\labelenumi{(\alph{enumi})}
\setcounter{enumi}{3}
\item
  证明：若 \(u, v\) 是 \(\mathbf{E} * \mathbf{F}\) 中使 \(h(u) = r\) , \(h(v) = s\) 的元素，则 \(h(uv) \leqslant h(u) + h(v)\) ，且 \(h(uv) = h(u) + h(v)\) ，除非 \(u\) 与 \(v\) 都是纯的，并且或者 \(r\) 是偶数而 \(u, v\) 奇偶性不同，或者 \(r\) 是奇数而 \(u, v\) 奇偶性相同。（例如设 \(r\) 是偶数， \(u \in \mathbf{P}_r'\) , \(v \in \mathbf{P}_s'\) ，证明 \(uv \in \mathbf{P}_{r+s-1}\) 仅当 \(u \in \mathbf{P}_{r-1}\) 或 \(v \in \mathbf{P}_{s-1}\) 时才可能。）
\item
  假设 E 没有零因子，且在 (d) 的记号下设 r 是偶数，u 是纯偶的而 v 是纯奇的。证明 \(h(uv) = r + s - 1\) ，除非存在 E 的两个可逆元素 x、y 使得 \(uy \in P_{r-1}\) 且 \(xb \in P_{s-1}\) 。（设 \((a_i)\) （相应地 \((b_j)\) ）是 (b) 中所考虑的 \(J_r + R_{r-1}\) 在 \(R_r\) 中（相应地 \(J_s + R_{s-1}\) 在 \(R_s\) 中）的补子空间的基；写出 \(u = \sum_i \phi(a_i)x_i\) 与 \(v = \sum_j y_j \phi(b_j)\) ，其中 \(x_i\) 与 \(y_j\) 都在 E 中，并证明：若 \(uv \in P_{r+s-2}\) ，则必然有 \(x_i y_j \in K\) 。）当设 E 与 F 都没有零因子时，类似地处理 \(h(uv) < h(u) + h(v)\) 的其他情形。
\item
  由(d)和(e)推出，当E和F没有零因子时，E * F也没有零因子。
\item
  将上述定义和结果推广到任意有限个含单位元的K-代数。
\end{enumerate}

\begin{enumerate}
\def\labelenumi{\arabic{enumi}.}
\setcounter{enumi}{6}
\tightlist
\item
  设 E 是域 K 上维数为 n 的向量空间。对每个整数 \(r \geqslant 1\) ，对称群 \(S_{r}\) 线性地作用于张量积 \(E^{\otimes r}\) 上，其作用为 \((\sigma, z) \mapsto \sigma \cdot z\) ，使得
\end{enumerate}

\[
\sigma . \left(x _ {1} \otimes x _ {2} \otimes \dots \otimes x _ {r}\right) = x _ {\sigma^ {- 1} (1)} \otimes x _ {\sigma^ {- 1} (2)} \otimes \dots \otimes x _ {\sigma^ {- 1} (r)}
\]

对于每个由 \(r\) 个向量组成的序列 \((x_i)_{1 \leqslant i \leqslant r}\) ，诸向量属于 \(\mathbf{E}\) 。设 \(u_{\sigma}(z) = \sigma, z\) 。证明：对每个序列 \((v_i)_{1 \leqslant i \leqslant r}\) （由 \(r\) 个 \(\mathbf{E}\) 的自同态组成），

\[
\operatorname{Tr} (u _ {\sigma} \circ (v _ {1} \otimes v _ {2} \otimes \dots \otimes v _ {r})) = \operatorname{Tr} (w _ {1}) \operatorname{Tr} (w _ {2}) \dots \operatorname{Tr} (w _ {k})\tag{*}
\]

其中 \(w_{j}\) 是 \(\mathbf{E}\) 的自同态，定义如下：把 \(\sigma^{-1}\) 分解为轮换， \(\sigma^{-1} = \lambda_1\lambda_2\ldots \lambda_k\) ，并且若 \(\lambda_{j} = (i_{1}i_{2}\dots i_{h})\) ，则记

\[
w _ {j} = v _ {i _ {1}} \circ v _ {i _ {2}} \circ \dots \circ v _ {i _ {h}}.
\]

（把它化为计算 (*) 的左端，其中每个 \(v_{i}\) 都形如 \(x \mapsto \langle x, a_{h}^{*} \rangle a_{k}\) ，这里 \((a_{j})_{1 \leqslant j \leqslant n}\) 是 E 的一个基，而 \((a_{j}^{*})_{1 \leqslant j \leqslant n}\) 是对偶基。）

¶ 8. 设 A 是一个整环，K 是其分式域；对于每个 A-模 M，令 τ(M) 表示其挠模（II，§ 7，第 10 号）。

\begin{enumerate}
\def\labelenumi{(\alph{enumi})}
\item
  证明：代数 \(\mathsf{T}(\mathbf{M})\) 没有零因子，当且仅当 A-模 \(\mathsf{T}(\mathbf{M})\) 是无挠的（注意 \(\mathsf{T}(\mathbf{M})_{(\mathbb{K})}\) 同构于 \(\mathsf{T}(\mathbf{M}_{(\mathbb{K})})\) ，并利用习题 5）。
\item
  设 \(u: \mathbf{M} \to \mathbf{M}'\) 是一个单射 A-模同态。证明 \(\operatorname{Ker}(\mathsf{T}(u)) \subset \tau(\mathsf{T}(\mathbf{M}))\) ；若 \(\mathbf{M}'\) 是无挠的，则 \(\operatorname{Ker}(\mathsf{T}(u)) = \tau(\mathsf{T}(\mathbf{M}))\) （参见 II，§ 7，第 10 号，命题 27）。
\item
  给出一个无挠 A-模 M 的例子，使得 \(\tau(\mathsf{T}(\mathbf{M}))\) 非零（II，§ 7，习题 31）。
\end{enumerate}

¶ 9. 设 A 是一个交换环，F 是整数集 I = {[}1, n{]}（其中 n ≥ 2）上的自由结合代数，而 X \(_{i}\) (1 ≤ i ≤ n) 是相应的未定元，于是 F 在 A 上的一个典范基由乘积 X \(_{i_1}\) X \(_{i_2}\) \ldots{} X \(_{i_r}\) 组成，其中 (i \(_{1}\) , \ldots， i \(_{r}\) ) 是 I 中元素的任意有限序列。

\begin{enumerate}
\def\labelenumi{(\alph{enumi})}
\item
  设 \(C_{2}\) 表示换位子 \([X_{i}, X_{j}] = X_{i}X_{j} - X_{j}X_{i}\) （\(1 \leqslant i < j \leqslant n\)）的集合。设 \(C_{m-1}\) 已有定义，\(C_{m}\) 表示元素 \([X_{i}, P] = X_{i}P - PX_{i}\) 的集合，其中 \(1 \leqslant i \leqslant n\) 而 P 取遍 \(C_{m-1}\) 。设 U 是 F 中由 l 与诸 \(C_{m}\) （\(m \geqslant 2\)）的并集生成的分次子代数。证明：对 \(1 \leqslant i \leqslant n\) ，有 \([X_{i}, P] \in U\) 对所有 \(P \in U\) （只需限于 P 是 \(\bigcup_{m} C_{m}\) 中元素之积的情形，并对因子个数作归纳论证）。
\item
  从现在起设 A 是域。对每个整数 \(m \geqslant 0\) ，设 \(\mathbf{U}_m\) 是 U 中次数为 \(m\) 的齐次元素所成的向量 A-空间（于是 \(\mathbf{U}_0 = \mathbf{A}\) , \(\mathbf{U}_1 = \{0\}\) ）；选取一组基 \(\mathbb{R}_{m,1}, \ldots, \mathbb{R}_{m,d_m}\) （在 \(\mathbf{U}_m\) 的每一个中），由 \(\bigcup_{k\leqslant m}\mathbf{C}_k\) 中元素的乘积组成。对每个有序整数对 \((m,k)\) （满足 \(m\geqslant 2,1\leqslant k\leqslant d_m\)），设 \(\mathbf{E}_{m,k}\) 是 U 中由诸 \(\mathbb{R}_{p,j}\) （满足 \(p > m\) 或 \(p = m\) 且 \(j\geqslant k\) ）生成的向量子 A-空间；证明 \(\mathbf{E}_{m,k}\) 是 U 的一个双边理想，且若 \(\mathbf{L}_{m,k}\) 是 F 中由 \(\mathbf{E}_{m,k},\mathbf{L}_{m,k}\) 生成的左理想，则它是 F 的一个双边理想。
\item
  对于每一个多重指标 \(\alpha = (\alpha_{1},\ldots ,\alpha_{n})\in \mathbf{N}^{n}\) ，我们记作
\end{enumerate}

\[
\mathrm{X} ^ {\alpha} = \mathrm{X} _ {1} ^ {\alpha_ {1}} \mathrm{X} _ {2} ^ {\alpha_ {2}} \dots \mathrm{X} _ {n} ^ {\alpha_ {n}}.
\]

证明元素 \(X^{\alpha}R_{m,k}\) ( \(m > 0, l \leqslant k \leqslant d_{m}\) ）构成 F 在 K 上的一个基。（为证明这些元素构成向量空间 F 的一个生成系，请证明每个元素 \(X_{i_{1}} \ldots X_{i_{r}}\) 都是它们的线性组合，对次数 r 作归纳。为看出 \(X^{\alpha}R_{m,k}\) 线性无关，用反证法论证：考虑其中一个 \(X_{i}\) ，并考察诸单项式 \(X^{\alpha}\) （它们出现在 \(X^{\alpha}R_{m,k}\) 之间的一个线性关系中）关于该变元的最大次数，再对此最大次数作归纳；为此，化归到 K 为无限的情形（必要时考虑 \(F \otimes_{K} K'\) ，其中 \(K'\) 是 K 的一个扩域），并注意到当 \(X_{i} + \lambda\) 代换某个 \(R_{m,k}\) 中的 \(X_{i}\) （对某个 \(\lambda \in K'\) ）时，元素 \(R_{m,k}\) 保持不变）。

\begin{enumerate}
\def\labelenumi{(\alph{enumi})}
\setcounter{enumi}{3}
\item
  证明：双边理想 \(\mathfrak{J}(\mathbf{F})\) ，即 \(\mathbf{F}\) 中由换位子 \([u,v] = uv - vu\) （u、v 为 \(\mathbf{F}\) 的元素）生成的理想，是诸 \(\mathbf{U}_m\) （\(m\geqslant 2\)）之和；\(\mathbf{F} / \mathfrak{J}(\mathbf{F})\) 同构于 \(\mathrm{A}[\mathrm{X}_1,\dots ,\mathrm{X}_n]\) ，从而是一个整环。
\item
  证明代数 U 不是自由结合代数，方法是证明商 \(\mathrm{U}/\Im(\mathrm{U})\) （其中 \(\Im(\mathrm{U})\) 由 U 中元素的换位子生成）含有幂零元素 \(\neq0\) 。（考虑 \(z'\) ，即 \(\mathrm{U}/\Im(\mathrm{U})\) 中元素 \(z = [\mathbf{X}_{i}, \mathbf{X}_{j}]\) 的像；\(z'^{2} \neq 0\) 但 \(z'^{4} = 0\) 。）
\end{enumerate}

\begin{enumerate}
\def\labelenumi{\arabic{enumi}.}
\setcounter{enumi}{9}
\tightlist
\item
  设 A 为一个交换环，
\end{enumerate}

\[
\mathrm{M} ^ {\prime} \xrightarrow {u} \mathrm{M} \xrightarrow {v} \mathrm{M} ^ {\prime \prime} \longrightarrow 0
\]

一个 A-模的正合序列，并设 \(n\) 为一个整数 \(>0\) 。设 \(L = \operatorname{Coker}(T^n(u))\) ，于是有正合序列

\[
\mathsf {T} ^ {n} (\mathbf {M} ^ {\prime}) \xrightarrow {\mathsf {T} ^ {n} (u)} \mathsf {T} ^ {n} (\mathbf {M}) \to \mathbf {L} \to 0.
\]

对集合 \(\{1, 2, \ldots, n\}\) 的每个子集 J ，我们记

\[
\mathrm{P} _ {J} = \mathrm{N} _ {1} \otimes \mathrm{N} _ {2} \otimes \dots \otimes \mathrm{N} _ {n}, \quad \text { where } \mathrm{N} _ {i} = \mathrm{M} ^ {\prime} \text { if } i \in \mathrm{J}, \mathrm{N} _ {i} = \mathrm{M} \text { if } i \notin \mathrm{J}
\]

\((\mathbf{P}_{\varnothing}\) 因此等于 \(\mathsf{T}^n (\mathbf{M}))\) 和

\[
\mathrm{Q} _ {\mathrm{J}} = \mathrm{N} _ {1} \otimes \mathrm{N} _ {2} \otimes \dots \otimes \mathrm{N} _ {n}, \quad \text { where } \mathrm{N} _ {i} = \mathrm{M} ^ {\prime} \text { if } i \in \mathrm{J}, \mathrm{N} _ {i} = \mathrm{M} ^ {\prime \prime} \text { if } i \notin \mathrm{J}
\]

\((Q_{\varnothing}\) 因此等于 \(T^{n}(M^{\prime \prime})\) 。）对每个整数 \(i\) （满足 \(0 \leqslant i \leqslant n\)），设 \(P^{(i)}\) 表示 \(T^{n}(M)\) 的子模，它由诸 \(\mathbf{P}_{\mathbf{J}}\) （\(\mathbf{J}\) 取遍满足 \(\operatorname{Card}(\mathbf{J}) = i\) 的一切子集）的像的并集生成；并设 \(\mathbf{L}^{(i)}\) 为 \(\mathbf{P}^{(i)}\) 在 \(\mathbf{L}\) 中的典范像。

\begin{enumerate}
\def\labelenumi{(\alph{enumi})}
\tightlist
\item
  证明存在一个典范的满同态
\end{enumerate}

\[
\bigoplus_ {\operatorname{Card} (J) = i} Q _ {J} \rightarrow L ^ {(i)} / L ^ {(i + 1)}
\]

(参见 II, § 5, 第 6 号, 命题 6)。

\begin{enumerate}
\def\labelenumi{(\alph{enumi})}
\setcounter{enumi}{1}
\tightlist
\item
  若序列 \(0 \to \mathbf{M}' \xrightarrow{u} \mathbf{M} \xrightarrow{v} \mathbf{M}'' \to 0\) 是正合的且分裂，证明在 (a) 中定义的同态是双射。
\end{enumerate}

